\chapter{K3 surfaces}\label{ch:k3}

\section{Definition and first examples}

\begin{definition}
A \emph{K3 surface} is a compact connected complex surface $X$ with trivial
canonical bundle, $\omega_X\cong\mathcal O_X$, and $H^1(X,\mathcal O_X)=0$.
\end{definition}

The name honours Kummer, K\"ahler and Kodaira, and the mountain K2. Every
K3 surface is K\"ahler \Lstatus~(Siu, \cite[Thm.~7.1.1]{Huybrechts2016}); the
projective ones are the algebraic K3 surfaces.

\begin{example}[quartics]
A smooth quartic surface $X\subset\PP^3$ is K3: adjunction gives
$\omega_X\cong\mathcal O_X(4-4)=\mathcal O_X$, and $H^1(\mathcal O_X)=0$ from
the exact sequence $0\to\mathcal O_{\PP^3}(-4)\to\mathcal O_{\PP^3}\to\mathcal O_X\to0$.
The Fermat quartic $x^4+y^4+z^4+w^4=0$ is the standard example; its Picard
number is $20$, the maximum (Chapter~\ref{ch:singular}).
\end{example}

\begin{example}[double planes]
The double cover of $\PP^2$ branched along a smooth sextic is K3. This is the
generic member of the two-dimensional family with $\NS\cong\lat2$.
\end{example}

\begin{example}[Kummer surfaces]\label{ex:kummer}
Let $A$ be a complex torus of dimension $2$ and $\iota=-1$. The quotient
$A/\iota$ has $16$ ordinary double points; its minimal resolution
$\Km(A)$ is a K3 surface, the \emph{Kummer surface} of $A$. The sixteen
exceptional curves are disjoint $(-2)$-curves whose classes span, together
with a certain half-sum, the \emph{Kummer lattice} $K\subset\NS(\Km A)$ of rank
$16$ and discriminant $2^6$ \Lstatus~\cite[Ch.~14.3]{Huybrechts2016}. Kummer
surfaces of products $E_1\times E_2$ of elliptic curves are the most
computable K3 surfaces of all, and return in Chapters~\ref{ch:elliptic}
and~\ref{ch:singular}.
\end{example}

\begin{example}[elliptic K3 surfaces]
A Weierstrass equation $y^2=x^3+a_4(t)x+a_6(t)$ with $\deg a_4\le8$,
$\deg a_6\le12$ (and not both too small) defines, after minimal resolution,
an elliptic K3 surface $X\to\PP^1$. Chapter~\ref{ch:elliptic} is devoted to
these.
\end{example}

\section{The Hodge diamond and the lattice $H^2$}

\begin{theorem}[\Lstatus]\label{thm:hodge}
For a K3 surface $X$: $b_1=0$, $b_2=22$, the Euler number is $e(X)=24$, and
the Hodge numbers are $h^{2,0}=h^{0,2}=1$, $h^{1,1}=20$. The intersection form
makes $H^2(X,\Z)$ an even unimodular lattice of signature $(3,19)$, hence
\[
  H^2(X,\Z)\;\cong\;\Lambda_{K3}=\U^3\oplus E_8(-1)^2
\]
by Theorem~\ref{thm:evenunimod} \cite[Ch.~1, Prop.~3.5]{Huybrechts2016},
\cite[VIII.3]{BHPV2004}.
\end{theorem}

Evenness is Wu's formula ($x^2\equiv\langle x,w_2\rangle$ and $w_2=c_1\bmod 2=0$),
unimodularity is Poincar\'e duality, and the signature is the Hodge index
theorem together with $b_2^+=2h^{2,0}+1=3$. All K3 surfaces thus share
\emph{one} lattice. The whole theory is about which line inside it is
$H^{2,0}$.

\section{N\'eron--Severi and transcendental lattices}

\begin{definition}
$\NS(X):=H^{1,1}(X)\cap H^2(X,\Z)$ is the \emph{N\'eron--Severi lattice}, a
primitive sublattice of $H^2(X,\Z)$; its rank $\rho(X)$ is the \emph{Picard
number}. The \emph{transcendental lattice} is $T(X):=\NS(X)^{\perp}$.
\end{definition}

For a K3 surface $\Pic(X)\cong\NS(X)$ since $H^1(\mathcal O_X)=0$, and by the
Lefschetz $(1,1)$ theorem $\NS(X)$ is exactly the set of integral classes
orthogonal to the holomorphic $2$-form:
\[
  \NS(X)=\{x\in H^2(X,\Z):\langle x,\sigma\rangle=0\},\qquad
  0\neq\sigma\in H^{2,0}(X).
\]
Hence $T(X)$ is the smallest primitive sublattice whose complexification
contains $\sigma$. Signatures: if $X$ is projective, $\NS(X)$ has signature
$(1,\rho-1)$ and $T(X)$ has signature $(2,20-\rho)$; in particular $\rho\le20$,
and since $T\otimes\C\supseteq H^{2,0}\oplus H^{0,2}$ we have $\rk T\ge2$.

\begin{theorem}[\Lstatus, minimality of $T$]\label{thm:Tminimal}
$T(X)_\Q$ is the smallest $\Q$-sub-Hodge structure of $H^2(X,\Q)$ containing
$H^{2,0}$; equivalently, $T(X)_\Q$ is an irreducible Hodge structure
\cite[Lem.~3.3.1]{Huybrechts2016}. If $X$ is projective, the Hodge structure
$T(X)_\Q$ is simple in the strong sense of Zarhin~\cite[Thm.~1.6]{Zarhin1983}:
its endomorphism algebra is a field, totally real or CM.
\end{theorem}

\begin{corollary}\label{cor:rho19}
If a one-parameter family of projective K3 surfaces has generic Picard number
$19$, then the generic transcendental lattice has rank $3$ and signature
$(2,1)$, and the Picard--Fuchs equation of a period has order $3$
(Chapter~\ref{ch:pf}).
\end{corollary}

\begin{warning}\label{warn:rho-from-arithmetic}
``$\rk T=22-\rho$'' is arithmetic, not evidence. If one \emph{assumes} a
family is $M_7$-polarized, then $\rho\ge19$ generically and $\rk T\le3$; the
statement $\rk T=3$ then says nothing beyond the assumption. A recurring
error in this subject is to present the number $3$ as corroboration of the
assumption that produced it. We shall flag the places where the temptation
arises.
\end{warning}

\section{Picard number $20$}

\begin{definition}
A K3 surface with $\rho(X)=20$ is called \emph{singular} (an unfortunate but
standard term; the surface is smooth). Then $T(X)$ is a positive-definite
even lattice of rank $2$, i.e.\ a positive-definite even binary quadratic
form $\bigl(\begin{smallmatrix}2a&b\\b&2c\end{smallmatrix}\bigr)$.
\end{definition}

\begin{theorem}[\Lstatus, Shioda--Inose]\label{thm:shioda-inose}
The map $X\mapsto T(X)$ is a bijection between singular K3 surfaces up to
isomorphism and positive-definite even binary forms up to $\SL_2(\Z)$
equivalence~\cite{ShiodaInose1977}. Every singular K3 surface is defined over
a number field, and its transcendental lattice determines it.
\end{theorem}

The forms $A_2=\bigl(\begin{smallmatrix}2&1\\1&2\end{smallmatrix}\bigr)$,
$\lat2\oplus\lat2$ and $\bigl(\begin{smallmatrix}2&1\\1&4\end{smallmatrix}\bigr)$
of Chapter~\ref{ch:lattices} are therefore the transcendental lattices of
three specific K3 surfaces: the first two are the Kummer-related surfaces
$X_3$ and $X_4$ of Shioda--Inose, of discriminants $3$ and $4$, the smallest
possible; the third has discriminant $7$. Chapter~\ref{ch:singular} explains
how they are constructed from products of elliptic curves with complex
multiplication.

\section{The ample cone and $(-2)$-curves}

For projective $X$, the positive cone
$\{x\in\NS(X)_\R:x^2>0\}$ has two components; call $\mathcal C_X$ the one
containing an ample class. The \emph{ample cone} is the set of $x\in\mathcal C_X$
with $\langle x,C\rangle>0$ for every $(-2)$-curve $C$ (a smooth rational
curve; $C^2=-2$ by adjunction). Equivalently \Lstatus~\cite[Cor.~8.1.6]{Huybrechts2016}:
the ample cone is the chamber of $\mathcal C_X$ cut out by the walls
$\delta^{\perp}$, $\delta\in\NS(X)$, $\delta^2=-2$, that contains an ample
class. This is the geometric meaning of the reflections and walls of
\S\ref{ex:swap}: the Weyl group of $\NS(X)$ acts on the chambers, and every
$(-2)$-class $\delta$ is, up to sign, effective (Riemann--Roch gives
$h^0(\delta)+h^0(-\delta)\ge2$).

\section{Exercises}

\begin{exercise}
Verify $e(X)=24$ for a smooth quartic from $c_2$ of the tangent bundle, and
for $\Km(A)$ from $e(A)=0$, the $16$ fixed points, and the $16$ exceptional
$\PP^1$'s.
\end{exercise}

\begin{exercise}[$\star$]
Let $\NS(X)\supseteq\U\oplus E_8(-1)^2\oplus\lat{-2n}=M_n$ primitively, with
$\rho=19$. Show that $T(X)\cong\U\oplus\lat{2n}$ using Example~\ref{ex:Mn},
and identify the exact hypothesis being used (primitivity, and equality of
$\NS$ with $M_n$).
\end{exercise}

\begin{exercise}
Show that a positive-definite even binary form of discriminant $3$ is
$\SL_2(\Z)$-equivalent to $A_2$, and one of discriminant $4$ to
$\lat2\oplus\lat2$. List the reduced forms of discriminants $7$, $28$, $40$.
\end{exercise}

\section*{Chapter note}

All theorems in this chapter are from Huybrechts~\cite{Huybrechts2016} and
BHPV~\cite{BHPV2004}, and Shioda--Inose~\cite{ShiodaInose1977}; theorem and
chapter numbers were checked against the cited editions to the extent the
author had them at hand, and are given so the reader can go straight to the
proof. The warning~\ref{warn:rho-from-arithmetic} records a lesson learnt in
the program that produced this book (Chapter~\ref{ch:frontier}).
